% P8 scorer result: artifact-backed XGBoost vs LLM, and the four operational reasons
\section{The Scorer Result: The Tree Keeps the Seat}
\label{sec:p8-scorer}

On the tabular scoring task, the current quick artifacts place XGBoost at test
AUC $0.7955$ on the $10{,}000$-row held-out split and the Qwen3.5-4B SFT
row-serialization arm at accuracy $0.792$ but AUC $0.48268$ on a 500-row
positive-enriched held-out evaluation (\tableref{tab:p8-headtohead}). The
accuracy is not evidence of useful ranking: at this class balance, AUC near
$0.5$ means the SFT scores do not order positives ahead of negatives. This is
consistent with the broader tabular literature: tree ensembles
remain the reference class on medium-sized tabular problems
\citep{grinsztajn2022tree, shwartzziv2021tabular}, and LLM row-serialization
approaches are most attractive in few-shot regimes, not at $40{,}000$ labeled
examples \citep{hegselmann2023tabllm}.

\begin{table}[t]
\centering
\caption{Current artifact-backed tabular scoring evidence on the custom
synthetic fraud dataset. XGBoost is evaluated on the full $10{,}000$-row test
split; the SFT LLM is evaluated on a 500-row positive-enriched held-out subset
because the natural-rate quick split has too few positives for stable LLM AUC.
Artifacts: \texttt{quick\_20260704/qp8\_fraud.tsv} and
\texttt{quick\_20260704/qp8-fraud-sft.tsv}. Operational columns are qualitative
rankings, argued in the text.}
\label{tab:p8-headtohead}
\small
\begin{tabular}{lcccccc}
\toprule
Model & Eval split & $n$ & Acc. & AUC $\uparrow$ & Latency & Injection surface \\
\midrule
XGBoost (scorer) & full test & $10{,}000$ & n/a & $\mathbf{0.7955}$ & ms-scale & none \\
Qwen3.5-4B SFT   & enriched test & $500$ & $0.792$ & $0.48268$ & s-scale & present \\
\bottomrule
\end{tabular}
\end{table}

The artifact-backed ranking is only the first of five reasons the tree keeps
the scorer seat. Even at AUC parity we would not move the LLM into the real-time scoring
path, for four operational reasons.

\paragraph{1. Latency.} Card authorization is a hard-real-time loop: the score
must arrive within a network-and-rules budget measured in tens of
milliseconds. A 200-tree, depth-6 ensemble scores a row in microseconds to
milliseconds on commodity CPUs; the committed run of our released script
measures roughly $6$\,ms per $10{,}000$ rows
(\texttt{xgboost\_results.json})---a rounding error against the authorization
budget. An LLM
forward pass over a serialized row---let alone an autoregressive
generation---is orders of magnitude slower and jitter-prone, and batching
tricks that amortize LLM cost are exactly what a per-authorization
synchronous path cannot use.

\paragraph{2. Cost.} The tree's marginal cost per transaction is effectively
zero at card-network volumes; the LLM's is a token bill that scales linearly
with traffic. Spending per-token inference on a task the tree does better is
the least defensible line item in the architecture we propose; the LLM budget
should be reserved for the low-volume, high-value tasks of
\secref{sec:p8-taxonomy}, where it has no substitute.

\paragraph{3. Calibration.} Downstream of the score sit threshold rules, alert
budgets, and expected-loss computations that consume the score \emph{as a
probability}. Boosted trees produce dense, monotonically usable scores that
calibrate well with standard post-hoc maps. LLM-verbalized confidences
inherit the miscalibration pathologies of modern neural networks
\citep{guo2017calibration}, compounded by the quantization of confidence into
tokens. A scorer whose $0.9$ does not mean ninety percent silently corrupts
every expected-loss decision behind it.

\paragraph{4. Prompt-injection exposure.} This reason is qualitative and, we
believe, decisive. A tree scores numbers; there is no channel through which a
transaction can \emph{address} the model. An LLM that reads serialized
transaction fields---merchant names, memo strings, cardholder-supplied
text---creates a new attack surface: adversarial instructions embedded in
attacker-controlled fields, the indirect prompt-injection pattern documented
by \citet{greshake2023injection}. A fraudster who learns that the scorer reads
free text will put text in front of it. Placing an instruction-following model
in a synchronous decision loop over attacker-authored inputs is a security
regression independent of accuracy, and it is the single strongest argument
for keeping the scorer seat non-linguistic.

The conclusion of this section is deliberately narrow: on fixed-schema tabular
scoring with ample labels, the current artifacts score the LLM below the
tree---and it would not get the seat even had it tied. Everything the LLM wins, it wins elsewhere in the pipeline---which is the
subject of the next section.
